\section{A criterion for the filtration}
\label{sec:assembly}

We first isolate the passage from compatible local bases to an actual
filtration. The construction removes the basis maps in blocks, ordered
by the largest index in their subsets. Each block kills the highest
remaining chromatic layer of the quotient. After the last positive
height has been removed, rational independence determines the remaining
block. Taking fibers of the quotient maps then recovers a filtration
of the original object. The same maps must serve as bases at every
height so that removing an earlier block has this prescribed effect
on every later test.

\subsection{Localizations and scalar conventions}

Fix a prime \(p\), and write \(S=S_p^\wedge\). We work in the stable
category of \(S\)-modules. Localizations are detected on the underlying
spectra. In particular, an \(S\)-module is \(E(r)\)-local or
\(K(t)\)-local when its underlying spectrum is so. These localizations
inherit their \(S\)-module structure and adjunction from the monoidal
Bousfield localization of spectra. All the fibers and cofibers below
are taken in \(S\)-modules.

We use three familiar consequences of chromatic localization:
\begin{enumerate}[label=(\roman*)]
\item If \(r\geq t\), the localization unit induces
\[
 L_{K(t)}S \xrightarrow{\;\simeq\;} L_{K(t)}L_rS.
\]
\item If an \(E(t)\)-local object \(Q\) satisfies \(L_{K(t)}Q=0\),
then \(Q\) is \(E(t-1)\)-local. This is the chromatic fracture square
applied to \(Q\).
\item \(L_0L_rS=L_0S=H\Qp\) for every \(r\geq0\).
Consequently rational \(S\)-modules are modules over \(H\Qp\).
\end{enumerate}
The first statement follows because every \(E(r)\)-equivalence is a
\(K(t)\)-equivalence. Composing localizations gives the first equality
in (iii): rationalization is already an \(E(r)\)-localization of lower
height. To identify its value, let \(\mathbb S\) denote the sphere before
completion. Its positive stable stems are finite: choose an odd sphere
in the stable range and apply Serre's finiteness theorem
\cite[Chapter~V, Section~3, Proposition~3]{Serre1951}.
The Moore exact sequences for \(\mathbb S/p^v\) and the Milnor sequence
for their inverse limit then give
\[
\pi_0S=\Zp,\qquad
\pi_iS=(\pi_i\mathbb S)^{\wedge}_p\quad(i>0),\qquad
\pi_iS=0\quad(i<0).
\]
Indeed, the finite Moore groups satisfy Mittag--Leffler, and the
transition on the torsion kernels in the Moore exact sequences is
multiplication by \(p\), so their inverse limit is zero. Thus the
positive groups displayed above are finite \(p\)-groups. Rationalizing
leaves only \(\Q\otimes\Zp=\Qp\) in degree zero, proving \(L_0S=H\Qp\).
The standard fracture square and these localization conventions are
reviewed in \cite[Section~2]{BB}.

We will also use that the unit of \(L_{K(1)}S\) is nonzero after
rationalization. At odd \(p\), height-one Morava theory is
\(E_1\simeq KU_p^\wedge\), with \(\Zp^\times\) acting through
Adams operations. The structured homotopy fixed point spectral sequence
\[
H_c^s(\Zp^\times;\pi_rE_1)
\Longrightarrow \pi_{r-s}L_{K(1)}S
\]
is strongly convergent
\cite[Theorem~1(iii)--(iv) and pp.~2--3]{DevinatzHopkins}.
The \(p\)-cohomological dimension of
\(\Zp^\times=\mu_{p-1}\times(1+p\Zp)\) is one. Since
\(\pi_*E_1\) is even and the action on \(\pi_0E_1=\Zp\) is
trivial, total degree zero has only the term
\(H_c^0(\Zp^\times;\pi_0E_1)=\Zp\). The augmentation sends the
actual sphere unit to \(1\) on this edge. Hence
\(\pi_0L_{K(1)}S\cong\Zp\), with the unit corresponding to \(1\);
its rationalization is nonzero.

\subsection{The simultaneous basis criterion}

For a finite subset \(I\) of the positive integers put
\[
d(I)=\sum_{i\in I}(2i-1),\qquad d(\varnothing)=0.
\]
When we refer to a map representing an element of
\(\pi_{-d(I)}D\), its source is the free \(S\)-module
\(\Sigma^{-d(I)}S\).

\begin{proposition}[Filtration criterion]
\label{prop:assembly}
Let \(n\geq1\), let \(p\) be odd, and let \(D\) be an \(S\)-module with
a specified map \(u:S\to D\). Suppose that maps
\[
z_I:\Sigma^{-d(I)}S\longrightarrow D,
\qquad I\subseteq\{1,\ldots,n-1\},
\qquad z_\varnothing=u,
\]
have the following properties.
\begin{enumerate}[label=(\alph*)]
\item For each \(1\leq t<n\), the map with the specified components
\begin{equation}
\bigvee_{I\subseteq\{1,\ldots,n-t\}}\Sigma^{-d(I)}S
 \xrightarrow{(z_I)} D
\label{eq:criterion-bases}
\end{equation}
is a \(K(t)\)-equivalence.
\item For every integer \(b\),
\begin{equation}
\dim_{\Qp}\pi_{-b}L_0D
=\#\{I\subseteq\{1,\ldots,n\}:d(I)=b\}.
\label{eq:criterion-rational}
\end{equation}
If \(n=1\), assume in addition that \(\pi_0L_0u\) is nonzero.
\end{enumerate}
Then \(X=L_{n-1}D\) has the filtration and unit compatibility of
Theorem~\ref{thm:main}(a), with \(u\) in place of its unit.
\end{proposition}

\begin{proof}
First suppose \(n>1\). We construct successive quotients \(Q_j\) of
\(X\), starting with \(Q_0=X\). The first block is
\[
B_1=L_{n-1}S\ \vee\ \Sigma^{-1}L_{n-1}S.
\]
Localizing the maps \(z_\varnothing,z_{\{1\}}\) gives a map
\(B_1\to Q_0\). Let \(Q_1\) be its cofiber. Applying
\(L_{K(n-1)}\), hypothesis~(a) identifies this map with an
equivalence. The cofiber \(Q_1\) is therefore \(E(n-2)\)-local.

Inductively, assume that \(2\leq j\leq n-1\) and that \(Q_{j-1}\)
is \(E(n-j)\)-local. For every \(I\) with \(\max I=j\), compose
\(z_I\) with \(D\to X\to Q_{j-1}\). The localization adjunction
extends this map uniquely, in the space of such extensions, to
\(\Sigma^{-d(I)}L_{n-j}S\). We thus obtain
\begin{equation}
B_j=\bigvee_{\max I=j}\Sigma^{-d(I)}L_{n-j}S
 \longrightarrow Q_{j-1}\longrightarrow Q_j,
\label{eq:block-quotient}
\end{equation}
where the last object is the cofiber.

Here is the compatibility needed to continue the induction. Fix
\(t=n-j\). Hypothesis~(a) gives a specified direct-sum basis of
\(L_{K(t)}D=L_{K(t)}X\), indexed by the subsets of
\(\{1,\ldots,j\}\). For each earlier block \(\ell<j\), its sphere
localization has height at least \(t\). Its \(K(t)\)-localization is
therefore the corresponding shift of \(L_{K(t)}S\), and its map is
the image of the original component \(z_I\). Successive cofibers
of these components remove exactly the subsets whose maximum is
less than \(j\), together with the empty subset. This assertion is
inductive: at each step the map into the cofiber is the composite
of that same original component with the quotient map, so the
cofiber is the quotient by that direct summand.

The remaining summands are precisely the components of \(B_j\).
It follows that \(L_{K(t)}(B_j\to Q_{j-1})\) is an equivalence.
Consequently \(L_{K(t)}Q_j=0\), and chromatic fracture makes
\(Q_j\) \(E(t-1)\)-local. The construction thus reaches a rational
object \(Q_{n-1}\). For \(n=2\), this is already the object \(Q_1\)
constructed from the first block.

We next determine that rational quotient, keeping track of possible
kernels. The elements
\[
z_I,\qquad I\subseteq\{1,\ldots,n-1\},
\]
are linearly independent after rationalization in each degree.
Indeed, under hypothesis~(a) at \(t=1\), their images in
\(L_{K(1)}D\) are the units of distinct shifted summands
\(\Sigma^{-d(I)}L_{K(1)}S\). The rationalized unit in each such
summand is nonzero. A rational relation between the \(z_I\) would
therefore give a relation between nonzero elements in different
summands, which is impossible. This argument includes subsets
with equal values of \(d(I)\).

Since \(L_0L_rS=H\Qp\), rationalizing the map from each block
gives, in every degree, the injection of the corresponding
independent vectors into the quotient by the preceding vectors.
The cofiber long exact sequence has no suspended kernel. Subtracting
these vectors from the dimensions in hypothesis~(b) leaves
\[
\dim_{\Qp}\pi_{-b}Q_{n-1}
=\#\{I\subseteq\{1,\ldots,n\}:\max I=n,\ d(I)=b\}.
\]
The derived category of \(\Qp\)-vector spaces splits by homology.
Choose a basis in each degree to obtain an equivalence of
rational \(S\)-modules
\[
B_n=\bigvee_{\max I=n}\Sigma^{-d(I)}H\Qp
 \xrightarrow{\;\simeq\;}Q_{n-1}.
\]
Its cofiber \(Q_n\) is zero.

It remains to pass from these successive quotients to a filtration
of \(X\). Set
\[
G_j=\fib(X\longrightarrow Q_j),\qquad 0\leq j\leq n.
\]
Then \(G_0=0\) and \(G_n=X\). For each block the square
\begin{equation}
\begin{tikzcd}
G_j \arrow[r] \arrow[d] & B_j \arrow[d]\\
X \arrow[r] & Q_{j-1}
\end{tikzcd}
\label{eq:filtration-pullback}
\end{equation}
is cartesian: its upper left corner is the fiber of
\(X\to Q_j=\cofib(B_j\to Q_{j-1})\). Taking fibers of its horizontal
maps yields the cofiber sequence
\[
G_{j-1}\longrightarrow G_j\longrightarrow B_j.
\]
If a block is \(B_j=C_1\vee\cdots\vee C_r\), replace \(G_j\) by
the intermediate pullbacks along
\[
0\longrightarrow C_1\longrightarrow C_1\vee C_2
\longrightarrow\cdots\longrightarrow B_j.
\]
Exactness of pullback in a stable category makes their successive
cofibers \(C_1,\ldots,C_r\). In the first block choose the unit
summand first. Its pullback identifies the first stage with
\(L_{n-1}S\), and its map to \(X\) is exactly the localized \(u\).
The other blocks are already ordered by maximum. The number of
individual pieces is
\[
2+\sum_{j=2}^{n-1}2^{j-1}+2^{n-1}=2^n.
\]
This constructs the required filtration by \(S\)-linear maps.

If \(n=1\), the object \(X=L_0D\) is rational. Hypothesis~(b)
places one copy of \(\Qp\) in degree zero and one in degree
\(-1\). The nonzero map \(L_0u:H\Qp\to X\) is an isomorphism on
degree-zero homotopy, so its cofiber is
\(\Sigma^{-1}H\Qp\). These are the two required stages.
\end{proof}

\begin{remark}
\label{rem:assembly-extensions}
The construction uses maps into the current quotient. Diagram
\eqref{eq:filtration-pullback} retains their boundary maps when
the blocks are pulled back to \(X\). Thus the criterion produces
the attaching maps as part of the filtration; it requires no
choice of a nullhomotopy for them.
\end{remark}
